The rapid adoption of Internet of Things (IoT) devices within domestic and small-to-medium business (SMB) environments has significantly expanded the digital attack surface, often without a corresponding increase in security awareness or technical capability among end users (Fortinet, 2024; Bugcrowd, 2026). Many consumer and SMB IoT deployments rely on devices that lack robust security controls, have limited visibility, and provide little or no notification when devices are compromised, taken offline, or interfered with through deliberate disruption such as Wi-Fi jamming (TechHive, 2024; Journal of Engineering 360, 2024).

This project presents the design and development of a security monitoring application that focuses on device availability, anomaly detection, and proactive alerting for IoT environments. The system is intended to operate without additional hardware, making it accessible and cost-effective for non-expert users.

The project adopts a full lifecycle approach, incorporating requirements analysis, system design, implementation, and evaluation. Core functionality includes device registration, continuous heartbeat monitoring, anomaly detection based on connectivity and behavioural indicators, and multi-channel alerting. The effectiveness of the proposed solution is evaluated through functional testing and simulated failure scenarios. The report concludes with a critical appraisal of the system's limitations and outlines potential future enhancements, including scalability improvements and advanced anomaly detection techniques.
